\chapter{Coûts de transaction dans la recherche}\label{ch:rs:transaction-costs-in-research}

Une prévision quotidienne portant sur cinquante grandes actions affiche un ratio de Sharpe brut de 2.56. Négociée comme le portefeuille moyenne--variance qu'elle implique, par un fonds de 1~milliard de dollars, elle fait tourner son capital
5.33 fois par jour et verse 1\,223~\% de son capital par an en coûts~: un ratio de Sharpe de $-29$ après coûts. La même prévision avec les coûts intégrés
dans l'optimiseur donne 1.80~; lissée sur une journée également, 2.43. Le signal n'a rien changé. Les coûts ne sont pas une décote appliquée à un \omterm{def:rs:vectorised-backtests:backtest}{backtest} achevé~; ils appartiennent
au modèle du marché, à l'optimiseur et au signal, et ils dépendent de la taille du fonds. Ce chapitre ajuste un modèle de coûts aux exécutions du cabinet, le place
dans la construction avec \texttt{firm.tcost}, compense entre elles les transactions des stratégies et rend les signaux sensibles aux coûts.

\section{Modèles de coûts}

\begin{definition}[Impact temporaire, impact permanent, loi d'impact en racine carrée, pourcentage du volume]\label{def:rs:transaction-costs-in-research:impact}
La part de la concession de prix d'une transaction qui s'estompe après celle-ci est l'\emph{impact temporaire}\index{impact temporaire}~; le \emph{impact
permanent}\index{impact permanent} est la part qui subsiste, l'information que le marché y lit. Le \emph{pourcentage du volume}\index{pourcentage du volume}
d'un ordre est sa taille $Q$ rapportée au volume quotidien $V$ du marché. La \emph{loi d'impact en racine carrée}\index{loi d'impact en racine carrée} énonce que l'impact moyen d'un ordre, en fraction
du prix, vaut environ $\eta\,\sigma\sqrt{Q/V}$ pour une volatilité quotidienne $\sigma$ et une constante $\eta$ d'ordre un.
\end{definition}

Le coût d'une transaction comporte trois parties~: la demi-fourchette (livre~1, chapitre~1), les frais et l'\omterm{def:rs:order-book-replay-simulation:impact}{impact de marché} (chapitre~18). Les deux premières sont linéaires en la taille échangée~;
l'impact ne l'est pas. Tóth et ses coauteurs ont expliqué la racine carrée par une offre et une demande qui s'annulent au prix courant, de sorte qu'un gros ordre doit être découpé et
digéré~; Almgren, Thum, Hauptmann et Li, qui ajustent près de 700\,000 ordres exécutés par les desks d'actions américains de Citigroup, trouvent que l'\omterm{def:rs:transaction-costs-in-research:impact}{impact temporaire} croît comme la puissance 3/5
du rythme de négociation. L'exposant voisin d'un demi compte plus que sa valeur précise~: une transaction deux fois plus grande coûte davantage par action, et le coût d'un portefeuille n'est donc pas
proportionnel à sa taille.

\section{Estimer les coûts à partir de ses propres exécutions}

Le modèle de coûts du cabinet vient de ses propres ordres parents. Le chapitre simule une année de ceux-ci, 5\,000 ordres de tailles allant de 0.01~\% à 20~\% du volume quotidien et
de volatilités quotidiennes allant de 1~\% à 3~\%, issus d'une loi implantée~: une demi-fourchette de 2 points de base et un impact $0.7\,\sigma\sqrt{Q/V}$, plus le mouvement propre du prix pendant que l'ordre
s'exécute. Le mouvement est la difficulté. À 1~\% du volume quotidien et 2~\% de volatilité, l'impact implanté vaut 14 points de base~; le mouvement du prix sur la durée de l'ordre a un
écart type de 63. Un ordre isolé ne dit presque rien de son coût, et un modèle de coûts est une régression sur des milliers d'ordres, dont le bruit croît avec la
taille de l'ordre.

\texttt{firm.tcost.fit\_impact} (le \cref{lst:rs:transaction-costs-in-research:fit}) moyenne les coûts dans vingt classes de participation, ajuste l'exposant sur les
logarithmes des moyennes par classe, et ajuste $\eta$ avec l'exposant fixé à un demi, avec une erreur type robuste au bruit croissant. À partir de 5\,000 ordres, il retrouve
$\eta = 0.71$ (erreur type 0.08) et un exposant de 0.48 (0.03)~; à partir de 500 ordres, $1.09$ (0.24) et 0.59 (0.07)~; à partir de 50\,000, $0.67$ (0.03) et 0.49 (0.01)
(la \cref{fig:rs:transaction-costs-in-research:fills}). Un desk ayant quelques centaines d'ordres par an connaît son coefficient de coût à un tiers près~; les hypothèses d'invariance de Kyle et Obizhaeva,
appuyées sur plus de 400\,000 ordres de transition de portefeuille, sont une façon d'emprunter
de la force aux autres titres en mettant les coûts à l'échelle du volume et de la volatilité de chaque titre.

\begin{omfigure}
\begin{tikzpicture}
\begin{loglogaxis}[width=10.5cm,height=6cm,xlabel={participation $Q/V$},ylabel={impact à 2~\% de volatilité (pb)},tick label style={font=\small},label style={font=\small},
  xmin=8e-5,xmax=0.25,ymin=0.5,ymax=150,grid=major,grid style={gray!25},legend pos=north west,legend style={font=\footnotesize},table/col sep=comma]
\addplot[only marks,mark=*,mark size=1.6pt,omDef,error bars/.cd,y dir=both,y explicit] table[x=participation,y=cost_bp,y error=se_bp]{figdata/research/27-transaction-costs-in-research/fills.csv};
\addplot[omThm,thick] table[x=participation,y=fit_bp]{figdata/research/27-transaction-costs-in-research/fills.csv};
\legend{moyennes par classe de 5\,000 ordres,loi en puissance ajustée}
\end{loglogaxis}
\end{tikzpicture}

\omcaption{Impact estimé à partir d'une année simulée d'ordres parents (demi-fourchette retirée, ramenée à une volatilité quotidienne de 2~\%), avec des barres d'erreur à deux erreurs types et la
loi ajustée. Les classes des plus petits ordres sont les plus bruitées par rapport à leur coût. Données~: \texttt{rs\_tcost.fills}.}\label{fig:rs:transaction-costs-in-research:fills}
\end{omfigure}

\section{Coûts dans l'optimiseur}

\begin{definition}[Optimisation tenant compte des coûts]\label{def:rs:transaction-costs-in-research:aware}
L'\emph{optimisation tenant compte des coûts}\index{optimisation tenant compte des coûts} maximise $\alpha^\top w - \frac\gamma2 w^\top\Sigma w - C(w - w_0)$, avec $C$ le modèle de coûts de la transaction
issu du portefeuille courant~: ici $C(d) = \sum_i s_i\lvert d_i\rvert + \eta\sigma_i\lvert d_i\rvert^{3/2}\sqrt{A/V_i}$ pour des demi-fourchettes $s_i$, un fonds de taille $A$ et des valeurs quotidiennes
échangées $V_i$.
\end{definition}

La puissance $3/2$ fait du problème un programme conique du second ordre (livre~4, chapitre~23). \texttt{firm.tcost.cost\_aware} le résout par une méthode de gradient proximal accéléré
plutôt~: la partie quadratique est lisse, le coût est séparable, et son pas proximal admet une forme fermée, un seuillage doux par la demi-fourchette suivi de la racine
d'un trinôme en $\sqrt{\lvert d\rvert}$ (le \cref{lst:rs:transaction-costs-in-research:aware}). Les portefeuilles du chapitre comptent les cinquante noms du chapitre~26, longs et courts, quotidiennement~; la
prévision est le rendement espéré véritable des trois signaux insérés, plus un bruit une fois et demie plus grand, renouvelé chaque jour, et $\gamma = 72$. Les valeurs quotidiennes
échangées de ces titres ont une médiane de 474~millions de dollars.

\begin{center}
\small
\begin{tabular}{@{}lrrrr@{}}
\toprule
taille du fonds & 10~M\$ & 100~M\$ & 1~Md\$ & 10~Md\$\\
\midrule
moyenne--variance, coûts ignorés~: après coûts & $-4.96$ & $-15.91$ & $-29.47$ & $-34.0$\\
coûts intégrés~: ratio de Sharpe avant coûts & 3.12 & 3.14 & 3.06 & 2.11\\
\quad après coûts & $-0.31$ & 0.06 & 1.80 & 1.76\\
\quad coûts annuels & 46.1~\% & 28.8~\% & 7.6~\% & 1.2~\%\\
\quad rotation quotidienne & 2.269 & 0.814 & 0.148 & 0.018\\
coûts intégrés, prévision lissée sur un jour & 1.68 & 2.03 & 2.43 & 1.66\\
coûts ignorés, prévision lissée sur cinquante jours & 1.78 & 1.63 & 1.15 & $-0.23$\\
\bottomrule
\end{tabular}
\end{center}

Le portefeuille qui ignore les coûts fait tourner son capital 5.33 fois par jour~; même un fonds de 10~millions de dollars, qui ne paie guère plus que la fourchette, perd. Le portefeuille tenant compte des coûts a
une surprise à la petite échelle~: il perd à 10~millions de dollars et gagne à 1~milliard de dollars. Avec des transactions bon marché, il croit à sa prévision et poursuit son bruit quotidien, que
l'optimiseur ne peut distinguer de l'alpha~; à 1~milliard de dollars, le coût de négociation le force à évoluer lentement, et évoluer lentement moyenne le bruit, de sorte que son ratio de Sharpe
avant coûts (3.06) est supérieur à celui du portefeuille naïf (2.56). Les coûts intégrés dans l'optimiseur ne protègent un portefeuille du bruit de sa prévision que lorsqu'ils sont grands par rapport à
ce bruit~; en deçà, la prévision elle-même doit être ralentie.

\section{Compensation entre stratégies}

\begin{definition}[Croisement interne]\label{def:rs:transaction-costs-in-research:crossing}
Le \emph{croisement interne}\index{croisement interne} compense les transactions des stratégies d'un même cabinet sur un même titre avant qu'aucune n'atteigne le marché~: un achat d'une stratégie et
une vente d'une autre se croisent sans coût de marché.
\end{definition}

Deux équipes négociant les mêmes cinquante titres à partir de leurs propres prévisions d'une même réalité, chacune avec son propre bruit, chacune lissée sur vingt jours et chacune gérant
1~milliard de dollars, paieraient 8.27~\% et 8.73~\% du capital d'une équipe par an, chacune de son côté. Compensées, leurs transactions coûtent 16.04~\%~: une économie de 0.96 point, 5.7~\% du total. L'économie
est plus faible que ne le laissent supposer les transactions opposées, car la compensation joue dans les deux sens sous une loi d'impact concave. Lorsque les équipes échangent en sens contraires, la transaction compensée est
plus petite et moins chère~; lorsqu'elles échangent dans le même sens, la transaction combinée est plus grande et, le coût croissant comme la puissance $3/2$, plus chère que les deux prises séparément~: 2.94 points par an
de plus ici. L'estimation de coût propre à chaque équipe sous-estime ce que ses transactions coûtent au cabinet lorsque d'autres négocient à ses côtés, ce qui annonce l'encombrement (crowding) traité au chapitre~28.
\texttt{firm.tcost.netting\_saving} répartit l'économie au prorata du coût autonome de chaque stratégie, de sorte que la compensation ne rende jamais
plus élevés les coûts déclarés d'une stratégie.

\section{Signaux sensibles aux coûts}

\begin{definition}[Lissage du signal, coût de point mort]\label{def:rs:transaction-costs-in-research:smoothing}
Le \emph{lissage du signal}\index{lissage du signal} remplace une prévision par une moyenne exponentiellement pondérée de ses valeurs récentes, avec une demi-vie choisie pour la
performance nette. Le \emph{coût de point mort}\index{coût de point mort} d'une stratégie est le coût par unité échangée pour lequel son rendement après coûts est nul~: son rendement brut
rapporté à sa rotation.
\end{definition}

Lisser un signal revient à échanger du signal contre de la rotation. Le lissage retarde la prévision, ce qui fait perdre la part de l'alpha qui se dégrade vite, et il moyenne le bruit de la prévision, ce qui
supprime des transactions qui n'avaient aucune raison d'être. La \Cref{fig:rs:transaction-costs-in-research:smoothing} montre le ratio de Sharpe net en fonction de la demi-vie, à 1~milliard de dollars. Le
portefeuille qui ignore les coûts passe de $-29.47$ à un maximum de 1.15 pour une demi-vie de cinquante jours, puis décroît au-delà (0.93 à cent jours, 0.69 à deux cents jours), car le retard coûte plus de
signal qu'il ne fait économiser de coûts. Le portefeuille qui tient compte des coûts atteint son maximum tôt, à 2.43 avec une demi-vie d'un jour, parce que l'optimiseur lisse déjà~; les demi-vies longues ne lui coûtent que du
signal (0.94 à deux cents jours). Le lissage optimal dépend de celui qui se charge du reste du lissage, et le \omterm{def:rs:portfolio-construction-ii:aim}{portefeuille cible} du chapitre~26 est la même idée rendue exacte
pour des signaux dont on connaît les taux de décroissance.

\begin{omfigure}
\begin{tikzpicture}
\begin{semilogxaxis}[width=11cm,height=6cm,xlabel={demi-vie du lissage (en jours~; le premier point correspond à l'absence de lissage)},ylabel={ratio de Sharpe après coûts},
  tick label style={font=\small},label style={font=\small},xmin=0.22,xmax=250,ymin=-4.2,ymax=3,xtick={0.3,1,10,100},xticklabels={aucun,1,10,100},
  grid=major,grid style={gray!25},legend pos=south east,legend style={font=\footnotesize},table/col sep=comma]
\addplot[black!60,thick,mark=o] table[x=half_life,y=naive]{figdata/research/27-transaction-costs-in-research/smoothing.csv};
\addplot[omDef,very thick,mark=*] table[x=half_life,y=aware]{figdata/research/27-transaction-costs-in-research/smoothing.csv};
\draw[black!50,densely dotted] (axis cs:0.22,0) -- (axis cs:250,0);
\legend{coûts ignorés,coûts intégrés à l'optimiseur}
\end{semilogxaxis}
\end{tikzpicture}

\omcaption{Le ratio de Sharpe net du portefeuille de cinquante titres à 1~milliard de dollars en fonction de la demi-vie du lissage de la prévision~; les valeurs inférieures à $-4$ sont tracées en $-4$. Données~:
\texttt{rs\_tcost.smoothing\_curve}.}\label{fig:rs:transaction-costs-in-research:smoothing}
\end{omfigure}

\section{Tutoriel~: de 2.56 à $-29$ et retour}\label{tut:rs:transaction-costs-in-research:tut}

\begin{tutorial}
\textbf{Objectif.} Ajuster la loi d'impact sur des exécutions simulées, négocier la prévision à cinquante titres en ignorant les coûts ou en les intégrant à l'optimiseur, pour quatre tailles de fonds, lisser la
prévision, compenser les transactions de deux équipes. \textbf{État final~:} le tableau, les \cref{fig:rs:transaction-costs-in-research:fills,fig:rs:transaction-costs-in-research:smoothing}.
\begin{tutsteps}
\item \textbf{L'ajustement}~: moyennes par classes pour l'exposant, ajustement robuste pour $\eta$.
\omcode{code/firm/tcost/firm_tcost.py}{34}{56}{Estimation de la loi d'impact à partir des exécutions simulées.}\label{lst:rs:transaction-costs-in-research:fit}
\item \textbf{Le portefeuille tenant compte des coûts}~: le pas proximal en forme close et l'itération accélérée.
\omcode{code/firm/tcost/firm_tcost.py}{67}{96}{Moyenne-variance avec fourchette et impact en loi de puissance intégrés.}\label{lst:rs:transaction-costs-in-research:aware}
\item \textbf{Exécuter} \texttt{rs\_tcost.estimate(n)}, \texttt{run(book, aum, half\_life)}, \texttt{smoothing\_curve}, \texttt{netting} et \texttt{fig\_tcost.py}.
\end{tutsteps}
\textbf{Pour aller plus loin.} Rendre le bruit de la prévision persistant (un processus autorégressif plutôt qu'une valeur nouvelle chaque jour) et observer l'allongement de la demi-vie optimale~; ajouter
seulement des coûts linéaires et voir apparaître la région de non-transaction du portefeuille tenant compte des coûts.
\end{tutorial}

\section{Construction~: la boîte à outils des coûts}\label{bld:rs:transaction-costs-in-research:tcost}

\begin{build}
\bfield{Rôle} Un seul modèle de coûts, estimé à partir des exécutions du cabinet et utilisé partout~: dans les \omterm{def:rs:vectorised-backtests:backtest}{backtests}, dans l'optimiseur, dans les estimations de capacité et dans la compensation
des stratégies.
\bfield{Interface} \texttt{impact\_bp(eta, sigma, participation, exponent)}, \texttt{fit\_impact(cost, sigma, participation, half\_spread)},
\texttt{trade\_cost(dw, aum, sigma, adv, half\_spread, eta, fee)}, \texttt{prox\_cost(v, s, c)},
\texttt{cost\_aware(alpha, Sigma, w0, aum, sigma, adv, half\_spread, eta, gamma, neutral)}, \texttt{net\_trades}, \texttt{netting\_saving(lists, cost\_fn)}, \texttt{smooth(signal, half\_life)}, \texttt{breakeven\_cost}.
\bfield{Règles} Les coûts dépendent de la taille du fonds et sont toujours imputés à la taille à laquelle la stratégie sera exploitée~; le modèle de coûts est réestimé sur les exécutions du cabinet, avec des erreurs
types~; aucune stratégie n'est jugée avant coûts~; les économies de compensation sont partagées de sorte qu'aucune stratégie ne voie ses coûts augmenter à cause d'une autre.
\bfield{Tests d'acceptation} \texttt{code/firm/tcost/tests/}~: la loi simulée retrouvée à l'intérieur de ses erreurs types~; le pas proximal optimal sur une grille fine~; le coût d'une
transaction, calculé à la main~; le portefeuille tenant compte des coûts égal à la moyenne-variance sans coûts, et localement optimal avec coûts~; la compensation calculée à la main~; la réponse indicielle du lissage~; le \omterm{def:rs:transaction-costs-in-research:smoothing}{coût de point mort}.
\bfield{Pour aller plus loin} \omterm{def:rs:transaction-costs-in-research:impact}{Impact temporaire} et permanent avec décroissance (un propagateur)~; coûts par plateforme et par moment de la journée~; construction multipériode tenant compte des coûts.
\end{build}

\begin{omsources}
\item B.~Tóth et al., ``Anomalous price impact and the critical nature of liquidity in financial markets'', \emph{Physical Review X} 1, 021006, 2011.
\item R.~Almgren, C.~Thum, E.~Hauptmann et H.~Li, ``Direct estimation of equity market impact'', \emph{Risk}, 2005.
\item A.~S.~Kyle et A.~A.~Obizhaeva, ``Market microstructure invariance: empirical hypotheses'', \emph{Econometrica} 84(4), 2016.
\item J.-P.~Bouchaud, J.~Bonart, J.~Donier et M.~Gould, \emph{Trades, Quotes and Prices}, Cambridge University Press, 2018.
\end{omsources}

\section{Exercices}

\begin{exercise}[$\star$]\label{exo:rs:transaction-costs-in-research:1}
Quel est le coût d'un ordre portant sur 5~\% du volume quotidien d'un titre dont la volatilité quotidienne est de 2~\%, avec une demi-fourchette de 2 points de base et $\eta = 0.7$~?
\end{exercise}

\begin{exercise}[$\star$]\label{exo:rs:transaction-costs-in-research:2}
Le portefeuille naïf rapporte 49.1~\% par an avant coûts et échange 5.33 fois son capital chaque jour. Quel est son \omterm{def:rs:transaction-costs-in-research:smoothing}{coût de point mort} par unité échangée~?
\end{exercise}

\begin{exercise}[$\star$]\label{exo:rs:transaction-costs-in-research:3}
Deux stratégies échangent chacune $q$ dans le même sens, sous un coût croissant comme $q^{3/2}$. De combien la transaction compensée coûte-t-elle plus cher que les deux transactions séparées~?
\end{exercise}

\begin{exercise}[$\star\star$]\label{exo:rs:transaction-costs-in-research:4}
Pourquoi un seul ordre parent en dit-il presque rien sur le coefficient de coût du cabinet~? Combien en faudrait-il~?
\end{exercise}

\begin{exercise}[$\star\star$]\label{exo:rs:transaction-costs-in-research:5}
Pourquoi le portefeuille tenant compte des coûts perd-il à 10~millions de dollars mais gagne-t-il à 1~milliard de dollars~?
\end{exercise}

\begin{exercise}[$\star\star$]\label{exo:rs:transaction-costs-in-research:6}
Établir le pas proximal de $s\lvert x\rvert + c\lvert x\rvert^{3/2}$.
\end{exercise}

\begin{exercise}[$\star\star\star$]\label{exo:rs:transaction-costs-in-research:7}
\emph{Programmation.} Exécuter \texttt{rs\_tcost.smoothing\_curve} pour les deux portefeuilles et expliquer pourquoi leurs demi-vies optimales diffèrent.
\end{exercise}

\begin{exercise}[$\star\star\star$]\label{exo:rs:transaction-costs-in-research:8}
\emph{Trouver la faille.} «~L'impact suit une racine carrée, donc nos coûts croissent comme la racine carrée de nos actifs~: doubler le fonds ne coûte que 41~\% de plus.~»
\end{exercise}

\section{Problème~: de 2.56 à $-29$}

\begin{problem}[{Problème du week-end --- un signal, chiffré}]\label{pb:rs:transaction-costs-in-research:1}
Les exécutions simulées du chapitre et les portefeuilles à cinquante titres.

\medskip\noindent\textbf{Partie I --- Le modèle de coûts.}
\begin{enumerate}
\item Énoncer la loi de coût simulée et l'ampleur du bruit de prix à 1~\% de participation.
\item Que retrouvent 500, 5\,000 et 50\,000 ordres pour $\eta$ et pour l'exposant~?
\item Que trouvent Almgren et ses coauteurs, et Tóth et ses coauteurs, pour l'exposant~?
\item Pourquoi l'ajustement se fait-il sur des moyennes par classes, et pourquoi une erreur type robuste~?
\end{enumerate}

\medskip\noindent\textbf{Partie II --- Les portefeuilles.}
\begin{enumerate}[resume]
\item Que valent la rotation, les coûts et le ratio de Sharpe net du portefeuille naïf à 1~milliard de dollars~?
\item Donner les ratios de Sharpe nets du portefeuille tenant compte des coûts pour les quatre tailles de fonds.
\item Pourquoi son ratio de Sharpe avant coûts est-il supérieur à celui du portefeuille naïf~?
\item Que se passe-t-il à 10~milliards de dollars~?
\end{enumerate}

\medskip\noindent\textbf{Partie III --- Netting et lissage.}
\begin{enumerate}[resume]
\item Que valent les coûts des deux équipes, séparés et après netting~?
\item Pourquoi le netting coûte-t-il davantage sur les transactions de même sens~?
\item Quels sont les ratios de Sharpe nets en fonction de la demi-vie pour les deux portefeuilles~?
\end{enumerate}

\medskip\noindent\textbf{Partie IV --- Le verdict.}
\begin{enumerate}[resume]
\item Énoncer le \emph{résultat nommé}~: le ratio de Sharpe net en fonction de la demi-vie du lissage, et le lissage optimal.
\item Lequel des deux portefeuilles faut-il faire tourner à 1~milliard de dollars~?
\item Quelle taille de fonds recommander pour le portefeuille naïf lissé~?
\item Comment un rapport de recherche doit-il présenter les coûts d'une stratégie~?
\item En quoi un modèle d'impact à propagateur changerait-il les conclusions~?
\item Comment estimer le modèle de coûts de la firme si elle ne dispose encore d'aucune exécution~?
\item Qu'implique le netting pour la capacité de plusieurs stratégies qui négocient les mêmes titres~?
\item Quel rapport le \omterm{def:rs:portfolio-construction-ii:aim}{portefeuille cible} du chapitre~26 entretient-il avec le lissage~?
\item En une phrase~: où les coûts doivent-ils figurer dans la recherche~?
\end{enumerate}
\end{problem}

\section{Questions d'entretien}

\begin{interviewq}[$\star$ \iqroles{trader, chercheur}]\label{iq:rs:transaction-costs-in-research:1}
Comment l'\omterm{def:rs:order-book-replay-simulation:impact}{impact de marché} évolue-t-il avec la taille de l'ordre~? Quels éléments empiriques le confirment~?
\end{interviewq}

\begin{interviewq}[$\star\star$ \iqroles{chercheur}]\label{iq:rs:transaction-costs-in-research:2}
Votre signal affiche un ratio de Sharpe quotidien de 2.5 avant coûts. Que faites-vous avant d'y croire~?
\end{interviewq}

\begin{interviewq}[$\star\star$ \iqroles{trader}]\label{iq:rs:transaction-costs-in-research:3}
Comment estimeriez-vous l'\omterm{def:rs:order-book-replay-simulation:impact}{impact de marché} de votre desk à partir de ses propres exécutions~?
\end{interviewq}

\begin{interviewq}[$\star\star$ \iqroles{chercheur, développeur}]\label{iq:rs:transaction-costs-in-research:4}
Comment intégrez-vous un coût en puissance $3/2$ dans un optimiseur de portefeuille~?
\end{interviewq}

\begin{interviewq}[$\star\star$ \iqroles{chercheur, trader}]\label{iq:rs:transaction-costs-in-research:5}
Deux stratégies de la firme négocient les mêmes actions. Que doit faire la firme à ce sujet~?
\end{interviewq}

\begin{interviewq}[$\star\star\star$ \iqroles{chercheur}]\label{iq:rs:transaction-costs-in-research:6}
Montrez que, lorsque les coûts croissent comme la puissance $3/2$ de la transaction, le rendement d'une stratégie après coûts est maximal pour une taille de fonds finie, et
la déterminer pour un rendement brut constant.
\end{interviewq}
